\section{Recovering a nilpotent operator from spaces of maps}
\label{p:probes}

Composition factors count the simple layers, and extensions record ways to
attach them. We now ask a reconstruction question with a complete answer:
can a suitable collection of maps into an object determine its isomorphism
class?

Our objects will be finite-dimensional vector spaces over an arbitrary field
$k$, equipped with a nilpotent operator. The answer will use only kernels,
images, quotients, and finite bases. In particular, we will prove the required
decomposition into chains. No classification of modules over a polynomial
ring, or Jordan normal form theorem, is assumed.

\subsection{A relation becomes a test object}
\label{p:relation-probe}

Let $V$ be a finite-dimensional $k$-vector space and let $N:V\to V$ be linear.
Recall how $N$ gives $V$ the structure of a $k[t]$-module $M$:
\[
 \left(\sum_i a_it^i\right)\!\cdot v=\sum_i a_iN^iv,
 \qquad N^0=\id_V.
\]
The identity $N^iN^j=N^{i+j}$ ensures that multiplication of polynomials
agrees with composition of their actions. A linear map $F:V\to V'$ is a
$k[t]$-module map precisely when $FN=N'F$: necessity follows by considering
$t$, and sufficiency follows by induction for powers of $t$ and then by
linearity for polynomials.

Assume henceforth that $N^R=0$ for some integer $R\geq1$. This is the
meaning of nilpotence. We allow $V=0$.

For every $r\geq1$, introduce the module
\[
 P_r=k[t]/(t^r).
\]
Its elements have unique representatives
$a_0+a_1t+\cdots+a_{r-1}t^{r-1}$: remove all terms divisible by $t^r$, and
observe that a polynomial of degree less than $r$ divisible by $t^r$ is zero.
Thus $1,t,\ldots,t^{r-1}$, understood as residue classes, form a $k$-basis.
Multiplication by $t$ runs along this basis and then gives zero.

The module $P_r$ has one generator with one relation, $t^r\cdot1=0$.
A map out of $P_r$ asks whether a proposed image of the generator satisfies
that relation.

\begin{proposition}[A kernel represented by a module]
\label{p:hom-kernel}
Evaluation at $1$ gives a $k$-linear isomorphism
\[
 \operatorname{ev}_1:\Hom_{k[t]}(P_r,M)\iso\Ker N^r,
 \qquad f\longmapsto f(1).
\]
Its inverse sends $v\in\Ker N^r$ to the map
\[
 f_v:P_r\longrightarrow M,\qquad \overline{p(t)}\longmapsto p(N)v.
\]
\end{proposition}

\begin{proof}
If $f$ is a module map, then
$N^rf(1)=f(t^r\cdot1)=0$. Moreover,
$f(\overline{p(t)})=p(N)f(1)$, so $f(1)$ determines $f$.

Conversely, if $N^rv=0$ and $p-p'=t^rq$, then
$(p-p')(N)v=q(N)N^rv=0$. Thus $f_v$ is well-defined on residue classes.
It is additive and $k$-linear, and
\[
 f_v(\overline{a(t)p(t)})=a(N)p(N)v
 =a(t)\cdot f_v(\overline{p(t)}),
\]
so it is a module map. The two constructions are inverse. Addition and
scalar multiplication of module maps are defined pointwise; the module
identities verify that these remain module maps. Evaluation and its inverse
respect these operations.
\end{proof}

This isomorphism respects maps between the objects being measured. If
$F:M\to M'$ is a module map, then $FN^r=(N')^rF$, so $F$ sends
$\Ker N^r$ into $\Ker(N')^r$. The square
\[
\begin{tikzcd}[column sep=large,row sep=large]
\Hom_{k[t]}(P_r,M) \arrow[r,"\operatorname{ev}_1","\sim"']
  \arrow[d,"f\mapsto F\circ f"']
  & \Ker N^r \arrow[d,"F|_{\Ker N^r}"] \\
\Hom_{k[t]}(P_r,M') \arrow[r,"\operatorname{ev}_1","\sim"']
  & \Ker (N')^r
\end{tikzcd}
\]
commutes because both paths send $f$ to $F(f(1))$. This compatibility is what
it means here for the identification to be \emph{natural} in $M$.

There is also compatibility as we enlarge the test object. The quotient
map $q_r:P_{r+1}\to P_r$ sends $\overline p$ to $\overline p$; it is
well-defined because $(t^{r+1})\subseteq(t^r)$. Precomposition by $q_r$
corresponds, under evaluation at $1$, to the inclusion
$\Ker N^r\subseteq\Ker N^{r+1}$. The growing sequence of kernels is thus
also a growing sequence of spaces of maps from explicit test objects.

\subsection{Which endpoints can be reached by a long chain?}
\label{p:bottom-spaces}

Set
\[
 K_r=\Ker N^r,\qquad d_r=\dim K_r\quad(r\geq0).
\]
Since $N^0=\id_V$, $K_0=0$ and $d_0=0$. For $r\geq R$, we have
$K_r=V$. The numbers $d_r$ are now well-defined by the dimension theorem
already proved. By Proposition~\ref{p:hom-kernel}, they are precisely the
dimensions of the spaces of maps from the $P_r$.

To understand these numbers, look at a second filtration, this time inside
$\Ker N$:
\[
 B_r=\Ker N\cap\im N^{r-1}\qquad(r\geq1).
\]
Every $B_r$ is a subspace, and
\[
 \Ker N=B_1\supseteq B_2\supseteq\cdots
 \supseteq B_R\supseteq B_{R+1}=0.
\]
An element of $B_r$ is killed by $N$ and can be reached by applying $N$
$r-1$ times. It can therefore be the endpoint of a chain of length $r$.
Membership in $B_{r+1}$ says that it can be reached by an even longer chain.

There is a short exact sequence of vector spaces
\begin{equation}
\label{p:kernel-bottom-ses}
 0\longrightarrow K_{r-1}\longrightarrow K_r
 \xrightarrow{\ N^{r-1}\ } B_r\longrightarrow0.
\end{equation}
The middle map has the stated target because $N^rx=0$ for $x\in K_r$.
Its kernel is $K_{r-1}$. It is surjective: if $y\in B_r$, write
$y=N^{r-1}x$; then $0=Ny=N^rx$, so $x\in K_r$. The first arrow is
inclusion, and the end arrows are the unique maps to or from the zero space.

Finite-dimensional dimension additivity applied to this sequence gives
\begin{equation}
\label{p:bottom-dim}
 \dim B_r=d_r-d_{r-1}.
\end{equation}
Consequently the successive increments in the sequence $(d_r)$ are
nonnegative and nonincreasing. Define
\begin{equation}
\label{p:canonical-multiplicity}
 E_r=B_r/B_{r+1},\qquad b_r=\dim E_r.
\end{equation}
Taking dimensions in the quotient gives
\begin{equation}
\label{p:second-difference}
 \boxed{\ b_r=2d_r-d_{r-1}-d_{r+1}.\ }
\end{equation}
This equation has been proved before any choice of chains. In particular,
its right side is nonnegative because it is the dimension of an actual
quotient space. We next prove why that quotient counts chains of exactly
length $r$.

\subsection{Lifting the endpoint filtration to a basis of chains}
\label{p:chain-construction}

For each $r$, choose vectors
\[
 u_{r,1},\ldots,u_{r,b_r}\in B_r
\]
whose residue classes form a basis of $B_r/B_{r+1}$. Such vectors exist by
finite-dimensional basis extension, proved after the dimension theorem.
Equivalently, start from the empty basis of $B_{R+1}=0$ and extend a basis
successively as we pass to $B_R,B_{R-1},\ldots,B_1$.

It follows that, for every $h\geq1$, the vectors
\begin{equation}
\label{p:endpoint-basis}
 \{u_{r,j}:h\leq r\leq R,\ 1\leq j\leq b_r\}
 \quad\text{form a basis of } B_h.
\end{equation}
Here is the elementary reason. If a subspace $W\subseteq U$ has a basis and
we lift a basis of $U/W$, the two lists together span $U$: subtract a
combination of the lifts to leave an element of $W$. They are independent:
reduce a relation modulo $W$ first, and then use independence in $W$.
Applying this observation down the filtration proves
\eqref{p:endpoint-basis}.

Because $u_{r,j}\in\im N^{r-1}$, choose $v_{r,j}\in V$ with
\[
 N^{r-1}v_{r,j}=u_{r,j}.
\]
Since $u_{r,j}\in\Ker N$, we also have $N^rv_{r,j}=0$. For
$1\leq h\leq r$, write
\[
 w_{r,j,h}=N^{r-h}v_{r,j}.
\]
The index $h$ is the height above zero: $w_{r,j,1}=u_{r,j}$,
$Nw_{r,j,1}=0$, and $Nw_{r,j,h}=w_{r,j,h-1}$ for $h>1$.
For example, a chain of length three has the form
\[
\begin{tikzcd}[column sep=large]
 w_{3,j,3} \arrow[r,mapsto,"N"] &
 w_{3,j,2} \arrow[r,mapsto,"N"] &
 w_{3,j,1} \arrow[r,mapsto,"N"] &0.
\end{tikzcd}
\]
The nodes in this last diagram are vectors, and the arrows show the action
of $N$ on them; it is not a diagram of vector spaces and linear maps.

\begin{theorem}[A basis of nilpotent chains]
\label{p:chain-basis}
The vectors $w_{r,j,h}$, for $1\leq r\leq R$, $1\leq j\leq b_r$ and
$1\leq h\leq r$, form a basis of $V$.
Consequently there is a $k[t]$-module isomorphism
\[
 M\cong\bigoplus_{r=1}^{R}P_r^{\oplus b_r}.
\]
\end{theorem}

\begin{proof}
First prove independence. Suppose that a linear combination of these
vectors is zero. If any coefficient is nonzero, let $H$ be the largest
height $h$ at which a nonzero coefficient occurs. Apply $N^{H-1}$ to the
relation. Every term of height less than $H$ vanishes, and a term
$w_{r,j,H}$ becomes $u_{r,j}$. We obtain a relation among the vectors
$u_{r,j}$ with $r\geq H$. These vectors are independent by
\eqref{p:endpoint-basis}. Thus every coefficient at height $H$ is zero,
contradicting its choice. No nonzero relation exists.

For spanning, prove by induction on $h$ that $K_h$ lies in the span of the
chosen chain vectors. The case $h=0$ is immediate. Given $x\in K_h$,
the vector $N^{h-1}x$ lies in $B_h$. By \eqref{p:endpoint-basis}, write
\[
 N^{h-1}x=\sum_{r\geq h,\,j}c_{r,j}u_{r,j}.
\]
Each $u_{r,j}$ in this sum has the chosen lift $w_{r,j,h}$ at height $h$.
Therefore
\[
 x-\sum_{r\geq h,\,j}c_{r,j}w_{r,j,h}\in K_{h-1}.
\]
The induction hypothesis puts this difference in the required span, and
so puts $x$ there. Since $K_R=V$, spanning follows.

For the module statement, on the copy of $P_r$ indexed by $(r,j)$ define
\[
 \overline{p(t)}\longmapsto p(N)v_{r,j}.
\]
This is a well-defined module map because $N^rv_{r,j}=0$, by the same
argument as Proposition~\ref{p:hom-kernel}. Combining these maps gives a map
from the displayed finite direct sum to $M$. It sends the standard basis
vectors $1,t,\ldots,t^{r-1}$ in each summand to the corresponding chain,
in reverse height order. The basis statement makes this map bijective;
its inverse respects the module operations because the map is bijective
and respects those operations. It is therefore a module isomorphism.
\end{proof}

The choices in this proof are visible. We selected representatives of
quotient bases and then lifted them through powers of $N$. Different
choices can give different chain bases. The spaces $B_r$, the quotients
$E_r$, and their dimensions $b_r$ depend only on $N$.

In particular, the construction proves the precise interpretation
\[
 \boxed{
 \dim\frac{\Ker N\cap\im N^{r-1}}
              {\Ker N\cap\im N^r}
 =\text{the number of chains of length exactly }r.
 }
\]
The quotient removes endpoints which can be lifted farther. A basis of
what remains supplies the new chains of length $r$ in the construction.

\subsection{Reading the operator from its measurements}
\label{p:profile-reconstruction}

In one chain of length $s$, the kernel of $N^r$ is spanned by the bottom
$\min(r,s)$ vectors. The direct-sum decomposition therefore gives
\begin{equation}
\label{p:profile-formula}
 d_r=\sum_{s=1}^{R}b_s\min(r,s),\qquad
 d_r-d_{r-1}=\sum_{s\geq r}b_s.
\end{equation}
The increment counts chains of length at least $r$. Subtracting successive
increments gives exactly the second-difference formula
\eqref{p:second-difference}, which we had already obtained from canonical
subquotients.

\begin{corollary}[Reconstruction by a finite nilpotent profile]
\label{p:reconstruction}
Two finite-dimensional nilpotent $k[t]$-modules $M$ and $M'$ are isomorphic
if and only if
\[
 \dim_k\Hom_{k[t]}(P_r,M)=\dim_k\Hom_{k[t]}(P_r,M')
 \qquad\text{for every }r\geq1.
\]
If a common nilpotence bound $R$ is known, it suffices to use $1\leq r\leq R$.
\end{corollary}

\begin{proof}
A module isomorphism identifies the spaces of maps by postcomposition,
so it preserves all the displayed dimensions. Conversely, equality of
the dimensions means equality of the $d_r$ by
Proposition~\ref{p:hom-kernel}. The formula
$b_r=2d_r-d_{r-1}-d_{r+1}$ then gives the same chain multiplicities for
both modules. Theorem~\ref{p:chain-basis} identifies each with the same
direct sum of the $P_r$, so they are isomorphic.

For a known common bound $R$, put $d_0=0$ and $d_{R+1}=d_R$; nilpotence
justifies the latter equality. The first $R$ measurements therefore
determine every possibly nonzero $b_r$.
\end{proof}

Thus the dimensions of these particular spaces of maps determine the
isomorphism class in this particular category. They do not choose an
isomorphism or a preferred basis. The proof succeeds because we established
a decomposition theorem and then showed that the measurements recover
every multiplicity in that decomposition.

Here are the first distinctions. Write $J_s$ for multiplication by $t$ on
$P_s$; as an operator, it is a single nilpotent chain of length $s$.
In dimension three we obtain
\begin{center}
\begin{tabular}{@{}lll@{}}
\toprule
Operator & $(d_1,d_2,d_3)$ & $(b_1,b_2,b_3)$ \\
\midrule
Zero on $k^3$ & $(3,3,3)$ & $(3,0,0)$ \\
$J_2\oplus J_1$ & $(2,3,3)$ & $(1,1,0)$ \\
$J_3$ & $(1,2,3)$ & $(0,0,1)$ \\
\bottomrule
\end{tabular}
\end{center}
The skipped-attachment operator from Section~\ref{x:skip-layer},
with $N(e_3)=e_1$ and $N(e_1)=N(e_2)=0$, is the middle row:
$(e_3,e_1)$ is a chain of length two and $e_2$ is a chain of length one.
Its first measurement is $d_1=2$, whereas the zero operator has $d_1=3$.
Thus even maps from $P_1$ recover the distinction that both adjacent
two-layer summaries missed.

All three have three simple factors isomorphic to
$S_0=P_1=k[t]/(t)$. Indeed, in a chain the subspaces spanned by its first
$h$ bottom vectors are submodules, and each successive quotient is a copy
of $S_0$. Concatenating these filtrations across the direct sum gives a
composition series. The space $S_0$ is simple because it is
one-dimensional over $k$ and every submodule is a $k$-subspace.
Consequently all three examples have the same class $3[S_0]$ in the
Grothendieck group. Their spaces of maps from $P_1,P_2,P_3$ distinguish them.

This does not contradict the universality of the additive invariant:
these dimensions of spaces of maps need not be additive on short exact
sequences. For example,
\[
 0\longrightarrow P_1\xrightarrow{\ 1\mapsto t\ }P_2
 \xrightarrow{\ \overline p\mapsto\overline p\ }P_1
 \longrightarrow0
\]
is exact. The first map identifies $P_1$ with the submodule $kt$ of
$P_2$, and the second map is reduction modulo $t$, with kernel $kt$.
Yet $\dim\Hom_{k[t]}(P_1,P_2)=1$, whereas the two endpoint measurements
sum to $1+1=2$. A module map $P_1\to P_2$ must send $1$ into $kt$, so its
composition with $P_2\to P_1$ is zero. In particular, the identity map
of the quotient $P_1$ cannot lift to $P_2$. The failure of additivity is
the lifting obstruction from the preceding chapter, appearing in an
explicit measurement.

For a slightly larger calculation, suppose measurements give
\[
 (d_0,d_1,d_2,d_3,d_4,d_5)=(0,3,5,6,7,7)
\]
for an operator with $N^4=0$. The increments are
$3,2,1,1,0$. Their successive differences give
\[
 (b_1,b_2,b_3,b_4)=(1,1,0,1).
\]
Hence $M\cong P_1\oplus P_2\oplus P_4$. The dimensions add to
$1+2+4=7$, and the kernel dimensions of this sum, calculated as
$\min(r,1)+\min(r,2)+\min(r,4)$, recover the original measurements.

The scope matters. These same tests do not reconstruct an arbitrary
operator. If $t$ acts as the identity on a nonzero vector space, every
$N^r$ is the identity and every $\Hom_{k[t]}(P_r,M)$ is zero. Its profile
is therefore the same as that of the zero module. The nilpotence
assumption ensured that our tests eventually reach all of $M$; without
it, an entire part of the object can remain undetected.

We have now proved a concrete reconstruction result. Composition factors
identify the simple layers. The canonical subquotients $E_r$ measure how
far their endpoints can be lifted along the operator. The test modules
$P_r$ turn those lifting depths into dimensions of spaces of maps, and the
chain-basis theorem shows that this information is sufficient to recover
the finite-dimensional nilpotent object up to isomorphism.
